\documentclass[11pt]{article}
\usepackage[a4paper,margin=18mm]{geometry}
\usepackage{fontspec}
\setmainfont{Latin Modern Roman}
\usepackage{amsmath,amssymb,amsthm,amscd,graphicx,color}
\usepackage{xurl}
\usepackage[unicode,hidelinks]{hyperref}
\allowdisplaybreaks
\setlength\emergencystretch{3em}

\makeatletter
\newtheorem{theorem}{Theorem}
\newtheorem*{theorem*}{Theorem}
\newtheorem{lemma}{Lemma}
\newtheorem*{lemma*}{Lemma}
\newtheorem{corollary}{Corollary}
\newtheorem*{corollary*}{Corollary}
\newtheorem{proposition}{Proposition}
\newtheorem*{proposition*}{Proposition}
\newtheorem{conjecture}{Conjecture}
\newtheorem*{conjecture*}{Conjecture}
{\theoremstyle{definition} \newtheorem{definition}{Definition}
\newtheorem*{definition*}{Definition}
\newtheorem{example}{Example}
\newtheorem*{example*}{Example}
\newtheorem{remark}{Remark}
\newtheorem*{remark*}{Remark}
\newtheorem{note}{Note}
\newtheorem*{note*}{Note}
}

\makeatother
\usepackage[all]{xy}
\usepackage{tikz}
\DeclareMathOperator{\co}{co}
\DeclareMathOperator{\Span}{Span}

\def\sw#1{{\sb{(#1)}}}

\newcommand{\ls}[1]{\ell(#1)^{\langle 1\rangle}}

\newcommand{\rs}[1]{\ell(#1)^{\langle 2\rangle}}

\newcommand{\lls}[2]{\ell_{#1}(#2)^{\langle 1\rangle}}

\newcommand{\rrs}[2]{\ell_{#1}(#2)^{\langle 2\rangle}}

\newcommand{\qcp}[1]{{C(\mathbb{R} P^{#1}_\mathcal{T})}}

\newcommand{\proj}[2]{\mathbb{#1}P_\mathcal{T}^#2}

\let\originalcedilla\c
\renewcommand{\c}[1]{{\fontspec{DejaVu Serif}\originalcedilla{#1}}}

\begin{document}
\begin{center}\Large Explicit strong connection on a finite piecewise principal right comodule algebra over a cocommutative Hopf algebra, with unital colinear local splittings\end{center}
\noindent\textbf{Stable ID:} 2023\par
\noindent\textbf{Authors:} Bartosz Zielinski\par
\noindent\textbf{Original work:} Piecewise Principal Coactions of Co-Commutative Hopf Algebras\par
\noindent\textbf{Version:} Official SIGMA 10 (2014) article 088, DOI 10.3842/SIGMA.2014.088; journal PDF and native TeX acquired 2026-10-01, exact bytes pinned by SHA256\par
\noindent\textbf{Selected task scope:} Full Theorem1 under its finite piecewise principal cocommutative-H and supplied-local-splitting hypotheses. The selected expression for ell uses ordered defect factors Ti from the source. Sections4-5 give constructive support for the local splittings; no assumption-free arbitrary Hopf-algebra generalization is asserted.\par
\noindent\textbf{Difficulty:} H2: The construction combines local strong connections with ordered convolution defect factors that telescope after applying every covering projection. The bundled multipullback reconstruction builds compatible colinear splittings by a double induction over the covering and a distributive-lattice basis construction, avoiding an assumption that a global strong connection has already been chosen.\par
\noindent\textbf{Editorial boundary note (not author text):} Full original Theorem 1 and its complete finite-piece strong-connection proof are retained with all piecewise-principality, covering and algebraic tensor conventions. H is cocommutative, every local connection lies in the specified V\_i tensor V\_i and each alpha\_i is unital and H-colinear on V\_i. The full new multipullback splitting induction, kernel-preserving gluing conditions, lattice basis partition and coinvariant lift construction are bundled as uncounted core, together with both colinearities and the convolution/telescoping counit identity. Earlier covering/principality and local strong-connection lifting criteria remain cited background. No separate splitting-helper count, Toeplitz-sphere example or noncocommutative extension is selected.\par
\noindent\textbf{Source:} \url{https://sigma-journal.com/2014/088/}\quad DOI: \url{https://doi.org/10.3842/SIGMA.2014.088}\par
\noindent\textbf{License and attribution:} Copyright retained by the named authors. Published by SIGMA under CC BY 4.0: \url{https://creativecommons.org/licenses/by/4.0/}. Source: \url{https://sigma-journal.com/about.html}. Adaptation consists of standalone excerpt formatting, cover and numbering restoration; original author language and mathematical text are retained.\par
\noindent\textbf{Editorial symbol notice (not author text):} Source Definition4 calls ideal sum/intersection meet/join in reversed lattice terminology; the actual distributivity and kernel conditions are explicit and the source bytes remain unchanged. Lemma3 calls B\_Gamma=\allowbreak{}pi(A\_Gamma) although its construction concerns basis vectors contained in the image; the preceding basis and lifting construction explicitly specify the selected representatives. No editorial rewriting of the proof.\par
\noindent\textbf{Typesetting adaptation (not author text):} The native cedilla accent in cited author names uses the portable DejaVu Serif font, because Latin Modern Roman lacks its combining cedilla glyph. Accent commands and bibliography text remain unchanged.\par
\medskip\hrule\medskip\noindent\textbf{Author text begins below.}\par
\setcounter{section}{1}
\setcounter{subsection}{0}
\setcounter{subsubsection}{0}
\setcounter{equation}{0}
% Original source lines 141--979
\section{Preliminaries}\label{preliminaries}

\subsection{Hopf algebra and comodule-related notation}

We work over a~f\/ixed ground f\/ield $\mathbb{K}$ and, unless stated otherwise, all vector spaces are understood to be
$\mathbb{K}$-vector spaces and the unadorned tensor product is understood to be the algebraic tensor product over~$\mathbb{K}$.
The comultiplication, counit and the antipode of a~Hopf algebra~$H$ are denoted by~$\Delta$, ${\epsilon}$ and~$S$,
respectively.
Let~$P$ be a~right comodule algebra.
We denote by $\Delta_P:P\rightarrow P\otimes H$ the right $H$-coaction on~$P$, and by
\begin{gather*}
P^{\co H}:=\big\{p\in P
\,|\,
\Delta_P(p)=p\otimes 1_H\big\}
\end{gather*}
the subalgebra of coaction invariant elements.
Instead of writing~$\Delta$'s and $\Delta_P$'s we usually employ the Heynemann--Sweedler notation with the summation
symbol suppressed, e.g.,
\begin{gather*}
\Delta(h)=:h\sw{1}\otimes h\sw{2},
\qquad
\Delta_P(p)=:p\sw{0}\otimes p\sw{1}.
\end{gather*}

\subsection{Principal comodule algebras}

Let~$H$ be a~Hopf algebra with bijective antipode, and let~$P$ be a~right $H$-comodule algebra.
Then~$P$ is a~principal comodule algebra if and only if there exists a~linear map
\begin{gather*}
\ell: \ H\rightarrow P\otimes P,
\qquad
\ell(h)=:\ls{h}\otimes\rs{h}
\end{gather*}
(note the Sweedler-like notation with summation sign supressed) satisfying the following conditions
\begin{subequations}
\label{el-prop}
\begin{gather}
\ell(1_H)=1_P\otimes1_P,
\label{el-unitality}
\\
\ls{h}\rs{h}={\epsilon}(h),
\label{el-collapse}
\\
\ls{h\sw{1}}\otimes\rs{h\sw{1}}\otimes h\sw{2}=\ls{h}\otimes\rs{h}\sw{0}\otimes\rs{h}\sw{1},
\label{el-rcolinear}
\\
S(h\sw{1})\otimes\ls{h\sw{2}}\otimes\rs{h\sw{2}}=\ls{h}\sw{1}\otimes\ls{h}\sw{0}\otimes\rs{h}.
\label{el-lcolinear}
\end{gather}
\end{subequations}

Such a~map, if it exists, is called a~{\em strong connection} on~$P$~\cite{bh04,dgh01, h-pm96}.
Strong connections are usually non-unique.

\subsection{Multi-pullbacks of algebras}

Let~$J$ be a~f\/inite set, and let
\begin{gather}
\label{family}
\big\{\pi^i_j:A_i\longrightarrow A_{ij}=A_{ji}\big\}_{i,j\in J,\,i\neq j}
\end{gather}
be a~family of algebra homomorphisms to which we will occasionally refer as ``gluing maps''.
\begin{definition}
[\cite{CM2000,GK-P99}] The \emph{multi-pullback algebra} $A^\pi$ of a~family~\eqref{family} of algebra homomorphisms is
def\/ined as
\begin{gather*}
A^\pi:=\bigg\{ (a_i)_{i\in J}\in\prod\limits_{i\in J}A_i
\,
\bigg|
\,
\pi^i_j(a_i)=\pi^j_i(a_j),
\;
\forall\, i,j\in J,\, i\neq j \bigg\}.
\end{gather*}
\end{definition}

\begin{definition}[\cite{cocycle}] A~family~\eqref{family} of algebra homomorphisms is called \emph{distributive} if and only if all of
them are surjective and their kernels generate distributive lattices of ideals.
\end{definition}

Let $(\pi^i_j:A_i\rightarrow A_{ij})_{i,j\in J,\,i\neq j}$ be a~family of surjective algebra homomorphisms.
For any distinct $i$, $j$, $k$ we put $A^i_{jk}:=A_i/(\ker\pi^i_j+\ker\pi^i_k)$ and take $[\cdot]^i_{jk}:A_i\rightarrow
A^i_{jk}$ to be the canonical surjections.
Next, we introduce the family of maps
\begin{gather*}
\pi^{ij}_k: \ A^i_{jk}\longrightarrow A_{ij}/\pi^i_j\big(\ker\pi^i_k\big),
\qquad
[a_i]^i_{jk}\longmapsto\pi^i_j(a_i)+\pi^i_j\big(\ker\pi^i_k\big).
\end{gather*}
They are isomorphisms when $\pi^i_j$'s are epimorphisms.

\begin{definition}
\label{cocycle-def}
We say \cite[Proposition~9]{CM2000} that a~family $(\pi^i_j:A_i\rightarrow A_{ij})_{i,j\in J,\,i\neq j}$ of algebra
epimorphisms satisf\/ies the {\em cocycle condition} if and only if, for all distinct $i,j,k\in J$,
\begin{enumerate}\itemsep=0pt
\item[1)] $\pi^i_j(\ker\pi^i_k)=\pi^j_i\big(\ker\pi^j_k\big)$,
\item[2)] the isomorphisms $\phi^{ij}_k:=\big(\pi^{ij}_k\big)^{-1}\circ\pi^{ji}_k:A^j_{ik}\rightarrow A^i_{jk}$ satisfy
$\phi^{ik}_j=\phi^{ij}_k\circ\phi^{jk}_i$.
\end{enumerate}
\end{definition}

Observe that, for all distinct $i,j,k\in J$ and any $a_i\in A_i$, $a_j\in A_j$,
\begin{gather}
\label{eqphi}
[a_i]^i_{jk}=\phi^{ij}_k\big([a_j]^j_{ik}\big)
\quad\!
\Leftrightarrow
\quad\!
\pi^{ji}_k\big([a_j]^j_{ik}\big)=\pi^{ij}_k\big([a_i]^{i}_{jk}\big)
\quad\!
\Leftrightarrow
\quad\!
\pi^i_j(a_i)-\pi^j_i(a_j)\in\pi^i_j\big(\ker\pi^i_k\big).\!\!\!
\end{gather}
One can prove (\cite{CM2000}, cf.~\cite{cocycle}, see also Theorem~\ref{main-split} in this paper) that the cocycle
condition together with distributivity guarantees that all projections on components of a~multipullback are surjective
(in fact all projections on submultipullbacks are surjective, but we will not make use of that fact).

\subsection{Piecewise principal comodule algebras}

\begin{definition}[\protect{cf.~\cite[Definition~3.7]{hkmz11}}] A~family of surjective algebra homomorphisms  $\{\pi_i:P\rightarrow P_i\}_{i\in\{1,\ldots,N\}}$
is called a~{\em covering}~\cite{hkmz11} if and only if
\begin{enumerate}\itemsep=0pt
\item[1)] $\bigcap_{i\in\{1,\ldots,N\}}\ker\pi_i=\{0\}$,
\item[2)] the family of ideals $(\ker\pi_i)_{i\in\{1,\ldots,N\}}$ generates a~distributive lattice with $+$ and $\cap$ as
meet and join, respectively.
\end{enumerate}
\end{definition}

Piecewise principal comodule algebras generalize the notion of (algebras of functions on) classical spaces which are
locally principal, but with respect to closed instead of open coverings~-- hence the use of the term ``piecewise''
instead of ``locally''.

\begin{definition}[\protect{see~\cite[Definition~3.8]{hkmz11}}]\label{piecewise-principal-def}
An $H$-comodule algebra~$P$ is called piecewise principal if there exists a~f\/inite
family $\{\pi_i:P\rightarrow P_i\}_{i\in J}$ of surjective $H$-comodule algebra morphisms such that
\begin{enumerate}\itemsep=0pt
\item[1)] the restrictions $\pi_i\big|_{P^{\co H}}:P^{\co H}\rightarrow P^{\co H}_i$ form a~covering,
\item[2)] the $P_i$'s are principal $H$-comodule algebras.
\end{enumerate}
\end{definition}

Note that, for all $i\in J$, $\pi_i(P^{\co H})\subseteq P^{\co H}_i$ by virtue of right $H$-colinearity of~$\pi_i$.
Hence, we were allowed to consider $\pi_i\big|_{P^{\co H}}$ in the statement of
Def\/inition~\ref{piecewise-principal-def} as a~map with codomain~$P^{\co H}_i$ without any additional assumptions.

By~\cite[Corollary~3.9]{hkmz11} a~piecewise principal comodule algebra is principal.
Note that any piecewise principal comodule algebra can be presented as a~multipullback comodule algebra with the gluing
maps being comodule algebra morphisms~\cite{cm2002}.

\section{Strong connection formula}\label{strong-section}

In this section we present an explicit (and arguably simple) expression for a~strong connection on a~piecewise
principal $H$-comodule algebra where~$H$ is a~co-commutative Hopf algebra.
Regretfully, the co-commutativity assumption is used crucially in the proof of the correctness of the formula, and so we
have little hopes of generalizing further the method which led to the derivation of this strong connection formula.

\begin{theorem}
\label{Main}
Let~$H$ be a~cocomutative Hopf algebra.
Let $\{\pi_i:P\rightarrow P_i\}_{i\in\{0,\ldots,n\}}$ be a~piecewise principal $H$-comodule algebra, and let
$\{\ell_i:H\rightarrow P_i\otimes P_i\}_{i\in\{0,\ldots,n\}}$ denote a~family of strong connections on $P_i$'s.
For any $i\in\{0,\ldots,n\}$, let $V_i$ be an~$H$ sub-comodule of $P_i$ such that $\ell_i(H)\subseteq V_i\otimes V_i$
and let $\alpha_i:V_i\rightarrow P$ be a~unital, colinear splitting of $\pi_i$, i.e.,
$\pi_i\circ\alpha_i=\mathrm{id}_{V_i}$.
For brevity, denote for $i\in\{0,\ldots,n\}$, $h\in H$
\begin{gather*}
\theta_i(h):={\epsilon}(h)-\alpha_i\big(\lls{i}{h}\big)\alpha_{i}\big(\rrs{i}{h}\big),
\\
T_i(h):=\theta_i(h\sw{1})\theta_{i+1}(h\sw{2})\cdots\theta_n(h\sw{n-i+1}),
\qquad
T_{n+1}(h):={\epsilon}(h).
\end{gather*}
Then the linear map $\ell:H\rightarrow P\otimes P$ defined for all $h\in H$ by the formula
\begin{gather*}
\ell(h)=\sum\limits_{i=0}^n\alpha_i\big(\lls{i}{h\sw{1}}\big)\otimes\alpha_i\big(\rrs{i}{h\sw{1}}\big)T_{i+1}(h\sw{2})
\end{gather*}
is a~strong connection on~$P$.
\end{theorem}

Note that, in particular, $T_n(h)=\theta_n(h)$, for all $h\in H$.
Note also that we consider splittings from $V_i$'s instead of splittings from $P_i$'s because the former are much easier
to construct.
\begin{proof}
Note that any co-commutative Hopf algebra has bijective (in fact involutive) antipode.
We need to prove that the map~$\ell$ def\/ined in the theorem, satisf\/ies all the properties~\eqref{el-prop}.

First note that, by the colinearity of $\alpha_j$'s, colinear properties~\eqref{el-lcolinear},~\eqref{el-rcolinear} of
$\ell_j$'s and the co-commutativity of~$H$ we have, that $\alpha_j(\lls{j}{h})\alpha_{j}(\rrs{j}{h})$ is a~coaction
invariant element of~$P$ for any $j\in\{0,\ldots,n\}$ and $h\in H$, and hence also $T_i(h)$ is a~coaction invariant
element of~$P$ for any $i\in\{0,\ldots,n+1\}$ and $h\in H$
\begin{gather*}
\rho^H\big(\alpha_j\big(\lls{j}{h}\big)\alpha_{j}\big(\rrs{j}{h}\big)\big)\\
\qquad{}
=\alpha_j\big(\lls{j}{h}\big)\sw{0}\alpha_{j}\big(\rrs{j}{h}\big)\sw{0} \otimes
\alpha_j\big(\lls{j}{h}\big)\sw{1}\alpha_{j}\big(\rrs{j}{h}\big)\sw{1}
\\
\qquad{}
 =\alpha_j\big(\lls{j}{h\sw{2}}\big)\alpha_{j}\big(\rrs{j}{h\sw{2}}\big)\otimes S(h\sw{1})h\sw{3}
\\
\qquad{}
 =\alpha_j\big(\lls{j}{h\sw{1}}\big)\alpha_{j}\big(\rrs{j}{h\sw{1}}\big)\otimes S(h\sw{2})h\sw{3}\\
 \qquad{}
 =\alpha_j\big(\lls{j}{h}\big)\alpha_{j}\big(\rrs{j}{h}\big)\otimes1.
\end{gather*}
In the penultimate equality we used co-commutativity of~$H$ to swap Sweedler indices $\sw{1}$ and $\sw{2}$ to be able to
use the antipode property.
In order to prove that~$\ell$ is left colinear (equation~\eqref{el-lcolinear}) we use the left colinearity of $\ell_i$'s
and the right colinearity of $\alpha_i$'s
\begin{gather*}
\ls{h}\sw{1}\otimes\ls{h}\sw{0}\otimes\rs{h}
\\
\qquad=
\sum\limits_{i=0}^n\alpha_i\big(\lls{i}{h\sw{1}}\big)\sw{1}\otimes\alpha_i\big(\lls{i}{h\sw{1}}\big)\sw{0}
\otimes\alpha_i\big(\rrs{i}{h\sw{1}}\big)T_{i+1}(h\sw{2})
\\
\qquad=
\sum\limits_{i=0}^n\lls{i}{h\sw{1}}\sw{1}\otimes\alpha_i(\lls{i}{h\sw{1}}\sw{0})
\otimes\alpha_i\big(\rrs{i}{h\sw{1}}\big)T_{i+1}(h\sw{2})
\\
\qquad=
\sum\limits_{i=0}^nS(h\sw{1})\otimes\alpha_i\big(\lls{i}{h\sw{2}}\big) \otimes\alpha_i\big(\rrs{i}{h\sw{2}}\big)T_{i+1}(h\sw{3})
\\
\qquad=
S(h\sw{1})\otimes\ls{h\sw{2}}\otimes\rs{h\sw{2}}.
\end{gather*}
The right colinearity %collinearity !!!
(equation~\eqref{el-rcolinear}) of~$\ell$ follows from the $H$-coaction invariance of $T_i(h)$'s,
the right colinearity of $\ell_i$'s, the right colinearity of $\alpha_i$'s, and the co-commutativity of~$H$
\begin{gather*}
\ls{h}\otimes\rs{h}\sw{0}\otimes\rs{h}\sw{1}
\\
\qquad=
\sum\limits_{i=0}^n\alpha_i\big(\lls{i}{h\sw{1}}\big)\otimes
\alpha_i\big(\rrs{i}{h\sw{1}}\big)\sw{0}T_{i+1}(h\sw{2})\otimes\alpha_i\big(\rrs{i}{h\sw{1}}\big)\sw{1}
\\
\qquad=
\sum\limits_{i=0}^n\alpha_i\big(\lls{i}{h\sw{1}}\big)\otimes
\alpha_i\big(\rrs{i}{h\sw{1}}\sw{0}\big)T_{i+1}(h\sw{2})\otimes\rrs{i}{h\sw{1}}\sw{1}
\\
\qquad=
\sum\limits_{i=0}^n\alpha_i\big(\lls{i}{h\sw{1}}\big)\otimes \alpha_i\big(\rrs{i}{h\sw{1}}\big)T_{i+1}(h\sw{3})\otimes h\sw{2}
\\
\qquad=
\sum\limits_{i=0}^n\alpha_i\big(\lls{i}{h\sw{1}}\big)\otimes \alpha_i\big(\rrs{i}{h\sw{1}}\big)T_{i+1}(h\sw{2})\otimes h\sw{3}
\\
\qquad=
\ls{h\sw{1}}\otimes\rs{h\sw{1}}\otimes h\sw{3}.
\end{gather*}
Here, in the penultimate inequality we used the co-commutativity of~$H$ exchanging Sweedler indices $\sw{2}$ and $\sw{3}$.

In order to prove that~$\ell$ is unital (equation~\eqref{el-unitality}), note f\/irst that for any $i\in\{0,\ldots,n\}$
\begin{gather*}
\theta_i(1)={\epsilon}(1)-\alpha_i\big(\lls{i}{1}\big)\alpha_{i}\big(\rrs{i}{1}\big)=1-\alpha_i(1)\alpha_i(1)=1-1=0,
\end{gather*}
because ${\epsilon}$, all $\ell_i$'s and all $\alpha_i$'s are unital.
It follows that $T_i(1)=0$ for all $i\in\{0,\ldots,n\}$, and $T_{n+1}={\epsilon}$ by def\/inition, hence
\begin{gather*}
\ell(1)=\sum\limits_{i=0}^n\!\alpha_i\big(\lls{i}{1}\big)\otimes\alpha_i\big(\rrs{i}{1}\big)T_{i+1}(1)
=\alpha_n\big(\lls{n}{1}\big)\otimes\alpha_n\big(\rrs{n}{1}\big)T_{n+1}(1)=1\otimes1,
\end{gather*}
where we used again the unitality of $\alpha_n$ and $\ell_n$.

Note now that for all $i\in\{0,\ldots,n\}$, and $h\in H$
\begin{gather*}
T_i(h)=T_{i+1}(h)-\alpha_i\big(\lls{i}{h\sw{1}}\big)\alpha_i\big(\rrs{i}{h\sw{1}}\big)T_{i+1}(h\sw{2}).
\end{gather*}
Indeed,
\begin{gather*}
T_i(h) =\theta_i(h\sw{1})T_{i+1}(h\sw{2})
 ={\epsilon}(h\sw{1})T_{i+1}(h\sw{2})-\alpha_i\big(\lls{i}{h\sw{1}}\big)\alpha_i\big(\rrs{i}{h\sw{1}}\big)T_{i+1}(h\sw{2})
\\
\phantom{T_i(h)}
 =T_{i+1}(h)-\alpha_i\big(\lls{i}{h\sw{1}}\big)\alpha_i\big(\rrs{i}{h\sw{1}}\big)T_{i+1}(h\sw{2}).
\end{gather*}
By applying this formula to $T_0(h)$ and keeping to expand with it the leftmost summand of the resulting expansion we
obtain easily
\begin{gather}
\label{t0}
T_0(h)={\epsilon}(h)-\sum\limits_{i=0}^n\alpha_i\big(\lls{i}{h\sw{1}}\big)\alpha_i\big(\rrs{i}{h\sw{1}}\big)T_{i+1}(h\sw{2}).
\end{gather}
On the other hand, for all $h\in H$ and $i\in\{0,\ldots,n\}$, as $\alpha_i$ is the splitting of $\pi_i$ it follows that
\begin{gather*}
\pi_i(\theta_i(h))  ={\epsilon}(h)-\pi_i\big(\alpha_i(\lls{i}{h})\big)\pi_i\big(\alpha_i(\rrs{i}{h})\big)
\\
\phantom{\pi_i(\theta_i(h))}
={\epsilon}(h)-\lls{i}{h}\rrs{i}{h} ={\epsilon}(h)-{\epsilon}(h)
 =0.
\end{gather*}
Hence
\begin{gather*}
\pi_i(T_j(h))=0,
\qquad
\text{for all}
\quad
i\geq j,
\quad
i\in\{0,\ldots,n\},
\quad
h\in H.
\end{gather*}
In particular, $\pi_i(T_0(h))=0$ for all $i\in\{0,\ldots,n\}$ and $h\in H$.
It follows that $T_0(h)=0$ for all $h\in H$ because $\bigcap_{i=0}^n\ker\pi_i=\{0\}$, as $\{\pi_i:P\rightarrow
P_i\}_{i\in\{0,\ldots,n\}}$ is a~covering.
The last fact is an immediate consequence of~\cite[Theorem~3.3]{hkmz11} and~\cite[Corollary~3.7]{hkmz11}.

Combining this with the equation~\eqref{t0} we obtain that for all $h\in H$
\begin{gather*}
\ls{h}\rs{h}=\sum\limits_{i=0}^n\alpha_i\big(\lls{i}{h\sw{1}}\big)\alpha_i\big(\rrs{i}{h\sw{1}}\big)T_{i+1}(h\sw{2})={\epsilon}(h),
\end{gather*}
i.e.,~$\ell$ satisf\/ies equation~\eqref{el-collapse} as needed.
\end{proof}

The expression for a~strong connection provided in the above theorem requires the unital and colinear splittings of
projections $\pi_i$ to be given.
The existence of such a~splittings is guaranteed by the~\cite[Lemma~3.1]{hkmz11} and~\cite[Theorem~3.3]{hkmz11}, but the
mere existence does not suf\/f\/ice for someone desirous of f\/inding the explicit formula.
The proof of~\cite[Lemma~3.1]{hkmz11} involves constructing a~unital and colinear splitting of surjective comodule
algebra map~$\pi$ from a~unital and linear splitting of restriction of~$\pi$ to the subalgebra of coaction invariant
elements (which always exists) utilizing the strong connection.
Hence, we cannot use even the slight simplif\/ication provided by the proof of~\cite[Lemma~3.1]{hkmz11}.

In practice, we expect that in many simpler cases, the appropriate splittings will not be dif\/f\/icult to guess.
However, for our result to be more widely applicable in practice, we will examine the explicit construction of colinear
and unital splittings of multipullback comodule algebra projections on components which does not assume the existence of
a~strong connection on a~multipullback comodule algebra (recall that a~piecewise principal comodule algebra can always
be presented as a~multipullback).

\section{Colinear splittings of piecewise principal comodule algebras}\label{colinear-section}

The result presented in this section allows to explicitly construct linear (colinear when approp\-riate) and unital
splittings of projections on components of a~multipullback (comodule) algebra.

\begin{theorem}
\label{main-split}
Suppose that a~family~\eqref{family} is distributive and satisfies the cocycle condition.
Moreover suppose that there exists two families $\alpha^i_j,\beta^i_j:A_{ij}\rightarrow A_i$, $i,j\in J$,
$j\neq i$ of linear $($colinear$)$ splittings of $\pi^i_j$'s such that all $\beta^i_j$'s are unital and for all distinct
$i,j,k\in J$ we have
\begin{gather}
\label{thkercond}
\alpha^i_j\big(\pi^i_j\big(\ker\pi^i_k\big)\big)\subseteq\ker\pi^i_k.
\end{gather}
Let $i\in J$, let $|J|=n+1$ and let $\kappa:\{0,\ldots,n\}\rightarrow J$ be a~bijection such that $\kappa_0=i$, where we
denote $\kappa_j:=\kappa(j)$ to easy the notation.
Then a~unital and linear $($colinear$)$ splitting $\alpha_i:A_i\rightarrow A^\pi$ of $\pi_i:A^\pi\rightarrow A_i$ can be
given explicitly, for any $a\in A_i$ as $\alpha_i(a):=(a_j)_{j\in J}$, where $a_i:=a$ and
$a_{\kappa_{m+1}}:=a^m_{\kappa_{m+1}}$ for any $0\leq m<n$.
The collections $\{a^k_{\kappa_{m+1}}\}_{0\leq k\leq m}\subseteq A_{\kappa_{m+1}}$, for $0\leq m<n$ are defined by the
following inductive formula
\begin{gather}
a^0_{\kappa_{m+1}}:=\beta^{\kappa_{m+1}}_{\kappa_0}\big(\pi^{\kappa_0}_{\kappa_{m+1}}(a_{\kappa_0})\big),
\nonumber
\\
a^{k+1}_{\kappa_{m+1}}:=a_{\kappa_{m+1}}^k-\alpha^{\kappa_{m+1}}_{\kappa_{k+1}}
\big(\pi^{\kappa_{m+1}}_{\kappa_{k+1}}\big(a^k_{\kappa_{m+1}}\big)-\pi^{\kappa_{k+1}}_{\kappa_{m+1}}(a_{\kappa_{k+1}})\big)
\label{recsplitform}
\end{gather}
for $0\leq k<m$.
\end{theorem}

\begin{proof}
It is clear that because all the maps involved in the def\/inition of $\alpha_i$ are unital and linear (colinear if need
be) then also $\alpha_i$ is linear (resp.
colinear).
The proof of unitality is slightly more subtle and it requires a~simple induction.
Pick some bijection $\kappa:\{0,\ldots,n\}\rightarrow J$ where $\kappa_0=i$.
Def\/ine $(a_j)_{j\in J}:=\alpha_i(1)$.
We need to show that $a_j=1$ for all $j\in J$.
Indeed, $a_{\kappa_0}=a_i=1$ by def\/inition.
Suppose we have proven that $a_j=1$ for all $0\leq j\leq m<n$.
Then using the equation~\eqref{recsplitform} we get
$a^0_{\kappa_{m+1}}=\beta^{\kappa_{m+1}}_{\kappa_0}(\pi^{\kappa_0}_{\kappa_{m+1}}(a_{\kappa_0}))=
\beta^{\kappa_{m+1}}_{\kappa_0}(\pi^{\kappa_0}_{\kappa_{m+1}}(1))=1$ as both~$\pi^{\kappa_0}_{\kappa_{m+1}}$ and~$\beta^{\kappa_{m+1}}_{\kappa_0}$ are unital.
Suppose now that we have proven that $a^k_{\kappa_{m+1}}=1$ for all $0\leq k<m$.
Then, equation~\eqref{recsplitform} yields
\begin{gather*}
a^{k+1}_{\kappa_{m+1}}
=a_{\kappa_{m+1}}^k-\alpha^{\kappa_{m+1}}_{\kappa_{k+1}}
\big(\pi^{\kappa_{m+1}}_{\kappa_{k+1}}(a^k_{\kappa_{m+1}})-\pi^{\kappa_{k+1}}_{\kappa_{m+1}}(a_{\kappa_{k+1}})\big)
\\
\phantom{a^{k+1}_{\kappa_{m+1}}}
 =1 -\alpha^{\kappa_{m+1}}_{\kappa_{k+1}}\big(\pi^{\kappa_{m+1}}_{\kappa_{k+1}}(1)-\pi^{\kappa_{k+1}}_{\kappa_{m+1}}(1)\big)
 =1-\alpha^{\kappa_{m+1}}_{\kappa_{k+1}}(0)
 =1.
\end{gather*}

Now it remains to show that $\alpha_i(a)\in A^\pi$ for all $a\in A_i$.
The inductive proof essentially follows the steps of the proof of~\cite[Proposition~9]{CM2000}.
We will show that for any $0\leq m\leq n$ we have
\begin{gather}
\label{partcond1}
\pi^{\kappa_j}_{\kappa_l}(a_{\kappa_j})=\pi^{\kappa_l}_{\kappa_j}(a_{\kappa_l}),
\qquad
\text{for all}
\quad
j,l\in\{0,\ldots,m\},
\quad
j\neq l.
\end{gather}
For $m=0$ this condition is emptily satisf\/ied.
Suppose we have proven the above condition for some~$m$.
In order to demonstrate it for $m+1$, we prove by induction that for any $0\leq k\leq m$, where $m<n$, we have
\begin{gather}
\label{partcond2}
\pi^{\kappa_j}_{\kappa_{m+1}}(a_{\kappa_j})=\pi^{\kappa_{m+1}}_{\kappa_j}\big(a^k_{\kappa_{m+1}}\big),
\qquad
\text{for all}
\quad
0\leq j\leq k.
\end{gather}
If $k=0$ then substituting the def\/inition of $a^0_{\kappa_{m+1}}$ yields (as $\beta^{\kappa_{m+1}}_{\kappa_0}$ is
a~splitting of $\pi^{\kappa_{m+1}}_{\kappa_0}$)
\begin{gather*}
\pi^{\kappa_{m+1}}_{\kappa_0}(a^0_{\kappa_{m+1}}) =\pi^{\kappa_{m+1}}_{\kappa_0}
\big(\beta^{\kappa_{m+1}}_{\kappa_0}\big(\pi^{\kappa_0}_{\kappa_{m+1}}(a_{\kappa_0})\big)\big)
 =\pi^{\kappa_0}_{\kappa_{m+1}}(a_{\kappa_0}).
\end{gather*}
Suppose now that we have proven condition~\eqref{partcond2} for some $0\leq k<m$.
Pick any $0\leq j\leq k$.
Then by (inductively assumed) condition~\eqref{partcond1} and equation~\eqref{eqphi} we have
\begin{gather}
\label{usedcoc}
[a_{\kappa_j}]^{\kappa_j}_{\kappa_{k+1}\kappa_{m+1}}=\phi^{\kappa_j\kappa_{k+1}}_{\kappa_{m+1}}
\big([a_{\kappa_{k+1}}]^{\kappa_{k+1}}_{\kappa_j\kappa_{m+1}}\big).
\end{gather}
Then it follows that
\begin{gather*}
\big[a^k_{\kappa_{m+1}}\big]^{\kappa_{m+1}}_{\kappa_j\kappa_{k+1}}
\overset{\substack{\text{by condition~\eqref{partcond2}} \\ \text{and equation~\eqref{eqphi}}}}{=}\phi^{\kappa_{m+1}\kappa_j}_{\kappa_{k+1}}
\big([a_{\kappa_j}]^{\kappa_j}_{\kappa_{m+1}\kappa_{k+1}}\big)
%\qquad \{\text{by condition~\eqref{partcond2} and equation~\eqref{eqphi}}\}
\\
\qquad{}
\overset{\text{by equation~\eqref{usedcoc}}}{=}
\phi^{\kappa_{m+1}\kappa_j}_{\kappa_{k+1}}\big(\phi^{\kappa_j\kappa_{k+1}}_{\kappa_{m+1}}
 \big([a_{\kappa_{k+1}}]^{\kappa_{k+1}}_{\kappa_j\kappa_{m+1}}
\big) \big)
%\qquad \{\text{by equation~\eqref{usedcoc}}\}
%\\
%\phantom{[a^k_{\kappa_{m+1}}]^{\kappa_{m+1}}_{\kappa_j\kappa_{k+1}}}{}
\overset{\substack{\text{by the cocycle} \\ \text{condition\vphantom{q}}}}{=} \phi^{\kappa_{m+1}\kappa_{k+1}}_{\kappa_j}\big([a_{\kappa_{k+1}}]^{\kappa_{k+1}}_{\kappa_j\kappa_{m+1}}\big).
%\qquad \{\text{by the cocycle condition}\}.
\end{gather*}
This equality, again by equation~\eqref{eqphi}, is equivalent to the following condition
\begin{gather*}
\pi^{\kappa_{m+1}}_{\kappa_{k+1}}\big(a^k_{\kappa_{m+1}}\big)-\pi^{\kappa_{k+1}}_{\kappa_{m+1}}(a_{\kappa_{k+1}})
\in\pi^{\kappa_{m+1}}_{\kappa_{k+1}}\big(\ker\pi^{\kappa_{m+1}}_{\kappa_j}\big).
\end{gather*}
Because the above relation ``is an element of'' holds for an arbitrary $0\leq j\leq k$ it implies immediately that
\begin{gather}
\label{kercond}
\pi^{\kappa_{m+1}}_{\kappa_{k+1}}\big(a^k_{\kappa_{m+1}}\big)-\pi^{\kappa_{k+1}}_{\kappa_{m+1}}(a_{\kappa_{k+1}})
\in\bigcap_{0\leq j\leq k}\pi^{\kappa_{m+1}}_{\kappa_{k+1}}\big(\ker\pi^{\kappa_{m+1}}_{\kappa_j}\big).
\end{gather}
Then
\begin{gather*}
\alpha^{\kappa_{m+1}}_{\kappa_{k+1}}\big(\pi^{\kappa_{m+1}}_{\kappa_{k+1}}\big(a^k_{\kappa_{m+1}}\big)
-\pi^{\kappa_{k+1}}_{\kappa_{m+1}}(a_{\kappa_{k+1}})\big)
\overset{\text{by condition~\eqref{kercond}}}{\in}
\alpha^{\kappa_{m+1}}_{\kappa_{k+1}}\left(\bigcap_{0\leq j\leq k}\pi^{\kappa_{m+1}}_{\kappa_{k+1}}\big(\ker\pi^{\kappa_{m+1}}_{\kappa_j}\big) \right)
%\qquad \{\text{By condition~\eqref{kercond}}\}  %\text
\\
\qquad
\overset{\text{by injectivity of $\alpha^{\kappa_{m+1}}_{\kappa_{k+1}}$}}{\in}
\bigcap_{0\leq j\leq k}
\alpha^{\kappa_{m+1}}_{\kappa_{k+1}}\left(\pi^{\kappa_{m+1}}_{\kappa_{k+1}}\big(\ker\pi^{\kappa_{m+1}}_{\kappa_j}\big)\right)
%\qquad \{\text{By injectivity of $\alpha^{\kappa_{m+1}}_{\kappa_{k+1}}$}\}
\overset{\text{by equation~\eqref{thkercond}}}{\subseteq}
\bigcap_{0\leq j\leq k}\ker\pi^{\kappa_{m+1}}_{\kappa_j},
%\qquad \{\text{By equation~\eqref{thkercond}}\},
\end{gather*}
that is
\begin{gather*}
%\label{diffin}
\alpha^{\kappa_{m+1}}_{\kappa_{k+1}}\big(\pi^{\kappa_{m+1}}_{\kappa_{k+1}}\big(a^k_{\kappa_{m+1}}\big)
-\pi^{\kappa_{k+1}}_{\kappa_{m+1}}(a_{\kappa_{k+1}})
\big)\in \bigcap_{0\leq j\leq k}\ker\pi^{\kappa_{m+1}}_{\kappa_j}.
\end{gather*}
The above equation implies immediately, that for all $0\leq l\leq k$
\begin{gather*}
\pi^{\kappa_{m+1}}_{\kappa_l}\big(a^{k+1}_{\kappa_{m+1}}\big)
= \pi^{\kappa_{m+1}}_{\kappa_l}\big(a_{\kappa_{m+1}}^k\big)-\pi^{\kappa_{m+1}}_{\kappa_l}\big(\alpha^{\kappa_{m+1}}_{\kappa_{k+1}}
\big(\pi^{\kappa_{m+1}}_{\kappa_{k+1}}\big(a^k_{\kappa_{m+1}}\big)-\pi^{\kappa_{k+1}}_{\kappa_{m+1}}(a_{\kappa_{k+1}})\big)\big)
\\
\phantom{\pi^{\kappa_{m+1}}_{\kappa_l}\big(a^{k+1}_{\kappa_{m+1}}\big)}{}
 =\pi^{\kappa_{m+1}}_{\kappa_l}\big(a_{\kappa_{m+1}}^k\big)
 =\pi^{\kappa_l}_{\kappa_{m+1}}(a_{\kappa_l}),
\end{gather*}
where, in the second equality we used the inductive assumption.
Moreover, using the fact that $\alpha^{\kappa_{m+1}}_{\kappa_{k+1}}$ is a~splitting of
$\pi^{\kappa_{m+1}}_{\kappa_{k+1}}$ we obtain
\begin{gather*}
\pi^{\kappa_{m+1}}_{\kappa_{k+1}}\big(a^{k+1}_{\kappa_{m+1}}\big)
= \pi^{\kappa_{m+1}}_{\kappa_{k+1}}\big(a_{\kappa_{m+1}}^k\big)-\pi^{\kappa_{m+1}}_{\kappa_{k+1}}\big(\alpha^{\kappa_{m+1}}_{\kappa_{k+1}}
\big(\pi^{\kappa_{m+1}}_{\kappa_{k+1}}\big(a^k_{\kappa_{m+1}}\big)-\pi^{\kappa_{k+1}}_{\kappa_{m+1}}(a_{\kappa_{k+1}})\big)\big)
\\
\phantom{\pi^{\kappa_{m+1}}_{\kappa_{k+1}}\big(a^{k+1}_{\kappa_{m+1}}\big)}
 =\pi^{\kappa_{m+1}}_{\kappa_{k+1}}\big(a_{\kappa_{m+1}}^k\big) -
\big(\pi^{\kappa_{m+1}}_{\kappa_{k+1}}\big(a^k_{\kappa_{m+1}}\big)-\pi^{\kappa_{k+1}}_{\kappa_{m+1}}(a_{\kappa_{k+1}})\big)
 =\pi^{\kappa_{k+1}}_{\kappa_{m+1}}(a_{\kappa_{k+1}}),
\end{gather*}
which ends the proof.
\end{proof}

At this point, the skeptical reader might be excused for doubting the applicability of Theorem~\ref{main-split}.
Indeed, while the existence of unital and linear splittings $\beta^i_j$'s of $\pi^i_j$'s follows immediately from the
surjectivity of $\pi^i_j$'s, and the existence of colinear splittings is assured (and assisted in explicit construction)
by~\cite[Lemma~3.1]{hkmz11} if all the $A_i$'s are principal comodule algebras, it is not clear how to f\/ind the linear
splittings $\alpha^i_j$ satisfying equation~\eqref{thkercond} nor that they exist at all in general case.
Fortunately, the results from the next section, interesting in their own right, not only assure the existence of
splittings $\alpha^i_j$ satisfying equation~\eqref{thkercond} under no stronger assumptions than those of
Theorem~\ref{main-split}, but they also provide the method of their (semi)-explicit construction.

\section{Colinear splittings of principal comodule algebras}\label{partitions-s}

\subsection{Partitions of sets}

Let~$A$ be a~set and let $A_i$, $i\in J$ be a~f\/ixed f\/inite family of subsets of~$A$.
For any $\Gamma\in 2^J$ we denote for brevity
\begin{gather}
\label{agamma}
A_\Gamma:=\bigcap_{i\in \Gamma}A_i.
\end{gather}
Obviously $A_{\Gamma_1}\cap A_{\Gamma_2}=A_{\Gamma_1\cup\Gamma_2}$.
Also $A_\varnothing = A$ by convention.
It is easy to see that $A_i$'s generate a~partition $\{B_\Gamma\}_{\Gamma\in 2^J}$ of~$A$ (i.e., all $B_\Gamma$'s are
disjoint and $A=\bigcup_{\Gamma\in2^J}B_\Gamma$) such that
\begin{gather*}
%\label{part-cond-set}
A_\Gamma=\bigcup_{\Gamma'\in2^J
\,| \,
\Gamma\subseteq\Gamma'}B_{\Gamma'},
\qquad
\text{for all}
\quad
\Gamma\in 2^J.
\end{gather*}
Indeed, the partition can be described explicitly, for all $\Gamma\in2^J$ by the formula
\begin{gather*}
B_\Gamma
:=\{x\in A
\,|\,
\forall\, i\in J,\,
x\in A_i\Leftrightarrow i\in \Gamma\}.
\end{gather*}

\subsection{Partitions of vector spaces}

Let now~$A$ be a~vector space and let $A_i$, $i\in J$ be a~f\/ixed f\/inite family of vector subspaces of~$A$.
$A_\Gamma$, for any $\Gamma\in 2^J$ is def\/ined as in equation~\eqref{agamma}.
We want to def\/ine a~linear counterpart of an associated partition $\{B_\Gamma\}_\Gamma$ def\/ined above for sets.
Similarly to plain sets, vector sub-spaces can be ordered by the set inclusion, and the resulting ordered set is
a~lattice, with subspace intersection ($V_1\cap V_2$) serving as inf\/imum and subspace sum ($V_1+V_2$) playing the role
of supremum.
The problem is that this lattice is not, in general, distributive.
It turns out that the assumption that the subspaces $A_i$, $i\in J$ generate a~distributive lattice is pivotal for
proving our desired result, stated immediately below:

\begin{lemma}
\label{partlem}
Let~$A$ be a~linear vector space and let $A_i$, $i\in I$ be a~finite family of vector subspaces of~$A$ generating a~distributive lattice.
$A$ has a~linear basis $\mathcal{B}=\bigcup_{\Gamma\in2^I}\mathcal{B}_\Gamma$, where
$\mathcal{B}_\Gamma\subseteq A_\Gamma$, $\Gamma\in2^I$, such that subsets $\mathcal{B}_\Gamma$ are all disjoint and
satisfy the following property
\begin{gather}
\label{partition}
A_\Gamma=\Span\left(\bigcup_{\Gamma'\in2^I,\; \Gamma'\supseteq\Gamma}\mathcal{B}_{\Gamma'}\right)  %\text  %\;
\end{gather}
for all $\Gamma\in2^I$.
\end{lemma}
\begin{proof}
First f\/ix a~linear order $\leq$ on $2^I$ subject to the condition
\begin{gather}
\label{orderdef}
\Gamma_1\supseteq \Gamma_2
\quad
\Rightarrow
\quad
\Gamma_1\leq\Gamma_2,
\qquad
\text{for all}
\quad
\Gamma_1,\Gamma_2\in 2^I.
\end{gather}
It is immediate that the minimal element in this order is~$I$ and maximal is $\varnothing$.
Note the following property of $\leq$ which will be used later
\begin{gather}
\label{Order}
\Gamma>\Gamma'
\quad
\Rightarrow
\quad
\Gamma\cup\Gamma'\supset \Gamma,
\qquad
\text{for all}
\quad
\Gamma,\Gamma'\in2^I.
\end{gather}
Indeed, assume $\Gamma>\Gamma'$.
$\Gamma\cup\Gamma'\supseteq \Gamma$ always, so we need just to show that the equality leads to contradiction.
Suppose that $\Gamma\cup\Gamma'= \Gamma$.
This is equivalent to $\Gamma\supseteq\Gamma'$ which implies by equation~\eqref{orderdef} that $\Gamma\leq\Gamma'$
contradicting the assumption $\Gamma>\Gamma'$.

The sets $\mathcal{B}_\Gamma$, $\Gamma\in2^I$ can be generated inductively (with respect to $\leq$) as follows
\begin{enumerate}\itemsep=0pt
\item[1)] $\mathcal{B}_I$ is some linear basis of $A_I$,
\item[2)] $\mathcal{B}_\Gamma$, for $\Gamma>I$, is chosen as a~maximal subset of $A_\Gamma$ such that
$\bigcup_{\Gamma'\leq\Gamma}\mathcal{B}_{\Gamma'}$ is linearly independent.
\end{enumerate}
It is immediate by construction of $\mathcal{B}_\Gamma$'s that $\mathcal{B}:=\bigcup_{\Gamma\in2^I}\mathcal{B}_\Gamma$
is a~linear basis of~$A$ and that all $\mathcal{B}_\Gamma$'s are disjoint.
Also by construction, $\mathcal{B}_{\Gamma'}\subseteq A_\Gamma$, $\Gamma\in2^I$ whenever $\Gamma\subseteq \Gamma'$,
which implies that half of property~\eqref{partition} is trivially satisf\/ied:
\begin{gather*}
%\label{partitionhalf}
\Span\left(\bigcup_{\Gamma'\in2^I,\, \Gamma'\supseteq\Gamma}\mathcal{B}_{\Gamma'}\right)\subseteq A_\Gamma
\end{gather*}
for all $\Gamma\in2^I$.
Finally, it is immediate that
\begin{gather}
\label{partition0}
A_\Gamma\subseteq\Span\left(\bigcup_{\Gamma'\in2^I,\, \Gamma'\leq\Gamma}\mathcal{B}_{\Gamma'}\right).
\end{gather}
We will prove the second half of property~\eqref{partition} by induction on $\leq$.

1.~$I$ is minimal in $2^I$ with respect to $\leq$.
Then by def\/inition of $B_I$ we have
\begin{gather*}
A_I= \Span(\mathcal{B}_I) = \Span\left(\bigcup_{\Gamma'\in2^I,\, \Gamma'\supseteq I}\mathcal{B}_{\Gamma'}\right).
\end{gather*}

2.~Suppose we have proven equation~\eqref{partition} for all $\Gamma'<\Gamma$.
For any $a\in A$, denote by $\{\alpha_\Gamma(a)\}_{\Gamma\in 2^I}$ the unique family of vectors such that
$a=\sum\limits_{\Gamma\in2^I}\alpha_\Gamma(a)$ and that $\alpha_\Gamma(a)\in \Span(\mathcal{B}_\Gamma)$ for all
$\Gamma\in2^I$ (they are unique because $\mathcal{B}$ is a~basis and $\mathcal{B}_\Gamma$'s are disjoint).
By~\eqref{partition0} $\alpha_{\Gamma'}(a)=0$ whenever $a\in A_\Gamma$ and $\Gamma'>\Gamma$, i.e.,
\begin{gather}
\label{partbase}
a=\sum\limits_{\Gamma'\in2^I,\, \Gamma'\leq\Gamma}\alpha_{\Gamma'}(a),
\qquad
\text{for all}
\quad
a\in A_\Gamma.
\end{gather}
Let $a\in A_\Gamma$.
Def\/ine $v:=a-\alpha_\Gamma(a)$.
By equation~\eqref{partbase}
\begin{gather*}
A_\Gamma\ni v = \sum\limits_{\Gamma'\in 2^I,
\,
\Gamma'<\Gamma}\alpha_{\Gamma'}(a)\in\sum\limits_{\Gamma'\in 2^I,
\,
\Gamma'<\Gamma}A_{\Gamma'},
\end{gather*}
hence
\begin{gather*}
v \in A_\Gamma\cap\left(\sum\limits_{\Gamma'\in 2^I,
\,
\Gamma'<\Gamma}A_{\Gamma'}\right)
\overset{\substack{\text{by distributivity of lattice}\\ \text{generated by $A_i$'s}}}{=}
\sum\limits_{\Gamma'\in 2^I,
\,
\Gamma'<\Gamma}A_{\Gamma'\cup\Gamma}
%\qquad \text{(by distributivity of lattice generated by $A_i$'s)}
\\
\qquad \overset{\substack{\text{by equation~\eqref{Order}} \\ \text{$\Gamma\subset\Gamma\cup\Gamma'$ if $\Gamma'<\Gamma$}}}{\subseteq}
\sum\limits_{\Gamma'\in 2^I,
\,
\Gamma'\supset\Gamma}A_{\Gamma'}
%\qquad (\text{by equation~\eqref{Order} $\Gamma\subset\Gamma\cup\Gamma'$ if $\Gamma'<\Gamma$})
 \overset{\substack{\text{by inductive assumption,} \\ \text{as $\Gamma'<\Gamma$ if $\Gamma'\supset\Gamma$}}}{\subseteq}
\Span\left(\bigcup_{\Gamma'\in 2^I,
\,
\Gamma'\supset\Gamma}\mathcal{B}_{\Gamma'} \right).
%\qquad (\text{by inductive assumption, as $\Gamma'<\Gamma$ if $\Gamma'\supset\Gamma$}).
\end{gather*}
It follows that
\begin{gather*}
a= \alpha_\Gamma(a)+v\in\Span(\mathcal{B}_\Gamma)+\Span\left(\bigcup_{\Gamma'\in 2^I,
\,
\Gamma'\supset\Gamma}\mathcal{B}_{\Gamma'} \right) = \Span\left(\bigcup_{\Gamma'\in 2^I,
\,
\Gamma'\supseteq\Gamma}\mathcal{B}_{\Gamma'} \right)
\end{gather*}
as needed.
\end{proof}

The following result is a~common knowledge:
\begin{lemma}
\label{common}
Let $\pi:A\rightarrow B$ be a~linear map, and let $\{A_i\}_{i\in I}$ be a~finite family of vector subspaces of~$A$.
Assume that $\ker\pi\cap\Big(\sum\limits_{i\in I}A_i\Big)=\sum\limits_{i\in I}(\ker\pi\cap A_i)$.
Then
\begin{gather*}
\pi\bigg(\bigcap_{i\in I}A_i\bigg)=\bigcap_{i\in I}\pi(A_i).
\end{gather*}
\end{lemma}

\begin{lemma}
\label{splitcoinv}
Let $\pi:A\rightarrow B$ be a~linear epimorphism, and let $\{A_i\}_{i\in I}$ be a~finite family of vector subspaces
of~$A$ such that $\{A_i\}_{i\in I}\cup\{\ker\pi\}$ generates a~distributive lattice of vector subspaces.
Then there exists a~linear splitting $\alpha:B\rightarrow A$ of~$\pi$ such that $\alpha(\pi(A_i))\subseteq A_i$ for all
$i\in I$.
\end{lemma}
\begin{proof}
Let $\mathcal{B}:=\bigcup_{\Gamma\in 2^I}\mathcal{B}_\Gamma$ be a~linear basis of~$B$ satisfying conditions guaranteed
by Lemma~\ref{partlem} with respect to the family $\{B_i\}_{i\in I}$, where $B_i:=\pi(A_i)$.
Note that Lemma~\ref{common} implies that $B_i$'s generate distributive lattice of ideals because $A_i$'s generate
distributive lattice of ideals.
For all $\Gamma\in2^I$ such that $\mathcal{B}_\Gamma$ is non-empty we def\/ine $\alpha(b)$ for all $b\in
\mathcal{B}_\Gamma$, to be an arbitrary element of $\pi^{-1}(b)\cap A_\Gamma$.
Note that $\pi^{-1}(b)\cap A_\Gamma$ is non-empty (so that this choice is possible) as $b\in B_\Gamma\neq \varnothing$,
and, $B_\Gamma=\pi(A_\Gamma)$ by Lemma~\ref{common}.
The map $\alpha:B\rightarrow A$ thus obtained is clearly a~linear splitting of~$\pi$.
For any $i\in I$ consider any $b\in B_i$.
Then, by Lemma~\ref{partlem}, $b\in\Span(\bigcup_{\Gamma\in 2^I  \, |
\,
i\in\Gamma}\mathcal{B}_\Gamma)$, and hence
\begin{gather*}
\alpha(b)\in\sum\limits_{\Gamma\in 2^I
\,
|
\,
i\in\Gamma}\sum\limits_{b'\in \mathcal{B}_\Gamma}\big(\pi^{-1}(b')\cap A_\Gamma\big) \subseteq\sum\limits_{\Gamma\in
2^I
\,
|
\,
i\in\Gamma}A_\Gamma\subseteq A_i.\tag*{\qed}
\end{gather*}
\renewcommand{\qed}{}
\end{proof}

Finally, we argue that we can generate a~colinear splitting (with appropriate properties) from the linear one on the
coaction invariant subalgebra:

\begin{lemma}
\label{alphaexists}
Let~$A$ be a~principal $H$-comodule algebra, let $\pi:A\rightarrow B$ be an $H$-comodule algebra surjection, and let
$\{A_i\}_{i\in I}$ be a~finite family of ideals in~$A$ which are subcomodules, such that $\{A_i\}_{i\in
I}\cup\{\ker\pi\}$ generates a~distributive lattice.
Define for all $i\in I$
\begin{gather*}
A_i^{\co H}:=A_i\cap A^{\co H},
\qquad
B_i:=\pi(A_i),
\qquad
B_i^{\co H}:=B^{\co H}\cap B_i.
\end{gather*}
Suppose that there exists a~linear map $\alpha^{\co H}:B^{\co H}\rightarrow A^{\co H}$ such that
\begin{gather*}
\pi\circ\alpha^{\co H}=\mathrm{id}_{B^{\co H}},
\qquad
\alpha^{\co H}\big(B^{\co H}_i\big)\subseteq A_i^{\co H},
\qquad
\text{for all}
\quad
i\in I.
\end{gather*}
Let $\ell:H\rightarrow A\otimes A$ be a~strong connection on~$A$.
Then the following formula
\begin{gather*}
\alpha: \ B\longrightarrow A,
\qquad
b\longmapsto \alpha^{\co H}\big(b\sw{0}\pi\big(\ls{b\sw{1}}\big)\big)\rs{b\sw{1}}
\end{gather*}
defines a~right $H$-colinear map satisfying
\begin{gather*}
\pi\circ\alpha=\mathrm{id}_{B},
\qquad
\alpha(B_i)\subseteq A_i,
\qquad
\text{for all}
\quad
i\in I.
\end{gather*}
\end{lemma}
\begin{proof}
The fact that~$\alpha$ def\/ined above is a~colinear splitting of~$\pi$ follows immediately from the proof of~\cite[Lemma~3.1]{hkmz11}.
It remains to show that $\alpha(B_i)\subseteq A_i$ for all $i\in I$.
Indeed, let $b\in B_i$.
Because of the left colinearity of~$\ell$ (equation~\eqref{el-lcolinear}) it follows easily that
$b\sw{0}\pi(\ls{b\sw{1}})\otimes\rs{b\sw{1}}\in B^{\co H}\otimes A$ (cf.\
proof of~\cite[Lemma~3.1]{hkmz11}), and because $B_i$ is an ideal and also a~right $H$-subcomodule, it follows also that
$b\sw{0}\pi(\ls{b\sw{1}})\otimes\rs{b\sw{1}}\in B_i\otimes A$, hence $b\sw{0}\pi(\ls{b\sw{1}})\otimes\rs{b\sw{1}}\in B^{\co H}_i\otimes A$.
Therefore
\begin{gather*}
\alpha(b)=\alpha^{\co H}\big(b\sw{0}\pi\big(\ls{b\sw{1}}\big)\big)\rs{b\sw{1}} \in \alpha^{\co H}\big(B_i^{\co
H}\big)A\subseteq A^{\co H}_iA\subseteq A_i.\tag*{\qed}
\end{gather*}
\renewcommand{\qed}{}
\end{proof}

Let us now put together all the steps needed to construct a~strong connection on a~piecewise principal comodule algebra
using the results presented in this paper.
Let~$H$ be a~co-commutative Hopf algebra.
Suppose that~$P$ is an $H$-comodule algebra which is piecewise principal with respect to a~f\/inite family
$\{\pi_i:P\rightarrow P_i\}_{i\in J}$ of surjective $H$-comodule algebra morphisms.
Then by~\cite[Corollary~3.7]{hkmz11} the family $\{\pi_i:P\rightarrow P_i\}_{i\in J}$ is a~covering.
Hence, by~\cite[Proposition~3]{CM2000}, it follows that~$P$ is isomorphic with the multipullback $P^\pi$, where the
gluing morphisms are def\/ined, for all $i,j\in J$, $i\neq j$,~by
\begin{gather}
\label{CanFamily}
\pi^i_j: \ P_i\longrightarrow P_{ij}:=P/(\ker\pi_i+\ker\pi_j),
\qquad
\pi_i(p)\longmapsto p+\ker\pi_i+\ker\pi_j.
\end{gather}
Obviously $\pi^i_j$'s are surjective.
Note that $\ker\pi^i_j=\pi_i(\ker\pi_j)$ for all $i,j\in J$, $i\neq j$.
Hence, as $\ker\pi_i$'s generate a~distributive lattice of ideals (see above), it follows by Lemma~\ref{common} that,
for all $i\in J$, also $\ker\pi^i_j$'s generate a~distributive lattice of ideals.
It is also immediate (see, e.g.,~\cite[Remark~2]{CM2000}) that the family~\eqref{CanFamily} satisf\/ies the cocycle
condition (Def\/inition~\ref{cocycle-def}).
Before we can use Theorem~\ref{main-split} to construct splittings of $\pi_i$'s needed by Theorem~\ref{Main} we still
need to construct splittings $\alpha^i_j,\beta^i_j:P_{ij}\rightarrow P_i$, $i,j\in J$,
$j\neq i$ of $\pi^i_j$'s with appropriate properties.
The $H$-colinear splittings $\beta^i_j:P_{ij}\rightarrow P_i$ can be constructed using~\cite[Lemma~3.1]{hkmz11} as all
the $P_{i}$'s are principal by assumption.
Similarly, using Lemma~\ref{alphaexists} we can construct $H$-colinear splittings $\alpha^i_j:P_{ij}\rightarrow P_i$ of
$\pi^i_j$'s satisfying $\alpha^i_j(\pi^i_j(\ker{}\pi^i_k)\subseteq\ker\pi^i_k$ for all distinct $i,j,k\in J$ from
a~family of linear maps $\big(\alpha^{\co H}\big)^i_j:P_i^{\co H}\rightarrow P_{ij}^{\co H}$, $i,j\in J$, $i\neq j$ satisfying
\begin{gather}
\label{props2alph}
\big(\pi^{\co H}\big)^i_j\circ\big(\alpha^{\co H}\big)^i_j=\mathrm{id}_{P_{ij}^{\co H}},
\qquad
\big(\alpha^{\co H}\big)^i_j\big(\big(\pi^{\co H}\big)^i_j\big(\ker\big(\pi^{\co H}\big)^i_k\big)\big)\subseteq\ker\big(\pi^{\co H}\big)^i_k,
\end{gather}
where we denoted by $\big(\pi^{\co H}\big)^i_j:P_i^{\co H}\rightarrow P_{ij}^{\co H}$, $i,j\in J$, $i\neq j$
the appropriate restrictions of $\pi^i_j$'s to the coaction invariant subalgebras.
Note that by~\cite[Lemma~3.1]{hkmz11} all the $(\pi^{\co H})^i_j$'s are surjective.
Then Lemma~\ref{splitcoinv} gives a~semi-explicit construction of $(\alpha^{\co H})^i_j$'s satisfying properties~\eqref{props2alph}.

Using $\alpha^i_j$'s and $\beta^i_j$'s we can now construct $H$-colinear and unital splittings of $\pi_i$'s using
Theo\-rem~\ref{main-split} and utilize them in the explicit construction of a~strong connection given~by
Theorem~\ref{Main}.
\begin{thebibliography}{99}
\bibitem[5]{bh04}
Brzezi{\'n}ski T., Hajac P.M., The {C}hern--{G}alois character,
  \href{http://dx.doi.org/10.1016/j.crma.2003.11.009}{\textit{C.~R.~Math. Acad. Sci. Paris}} \textbf{338} (2004), 113--116,
  \href{http://arxiv.org/abs/math.KT/0306436}{math.KT/0306436}.

\bibitem[6]{CM2000}
Calow D., Matthes R., Covering and gluing of algebras and differential
  algebras, \href{http://dx.doi.org/10.1016/S0393-0440(99)00035-2}{\textit{J.~Geom. Phys.}} \textbf{32} (2000), 364--396,
  \href{http://arxiv.org/abs/math.QA/9910031}{math.QA/9910031}.

\bibitem[7]{cm2002}
Calow D., Matthes R., Connections on locally trivial quantum principal fibre
  bundles, \href{http://dx.doi.org/10.1016/S0393-0440(01)00050-X}{\textit{J.~Geom. Phys.}} \textbf{41} (2002), 114--165,
  \href{http://arxiv.org/abs/math.QA/0002228}{math.QA/0002228}.

\bibitem[9]{dgh01}
D\c{a}browski L., Grosse H., Hajac P.M., Strong connections and
  {C}hern--{C}onnes pairing in the {H}opf--{G}alois theory, \href{http://dx.doi.org/10.1007/s002200100433}{\textit{Comm. Math.
  Phys.}} \textbf{220} (2001), 301--331, \href{http://arxiv.org/abs/math.QA/9912239}{math.QA/9912239}.

\bibitem[11]{h-pm96}
Hajac P.M., Strong connections on quantum principal bundles, \href{http://dx.doi.org/10.1007/BF02506418}{\textit{Comm.
  Math. Phys.}} \textbf{182} (1996), 579--617, \href{http://arxiv.org/abs/hep-th/9406129}{hep-th/9406129}.

\bibitem[12]{hkmz11}
Hajac P.M., Kr{\"a}hmer U., Matthes R., Zieli{\'n}ski B., Piecewise principal
  comodule algebras, \href{http://dx.doi.org/10.4171/JNCG/88}{\textit{J.~Noncommut. Geom.}} \textbf{5} (2011), 591--614,
  \href{http://arxiv.org/abs/0707.1344}{arXiv:0707.1344}.

\bibitem[19]{cocycle}
Hajac P.M., Zieli{\'n}ski B., Cocycle condition for multi-pullbacks of
  algebras, in Operator Algebras and Quantum Groups, \href{http://dx.doi.org/10.4064/bc98-0-9}{\textit{Banach Center
  Publ.}}, Vol.~98, Polish Acad. Sci., Warsaw, 2012, 239--243,
  \href{http://arxiv.org/abs/1207.0087}{arXiv:1207.0087}.

\bibitem[20]{GK-P99}
Pedersen G.K., Pullback and pushout constructions in {$C^*$}-algebra theory,
  \href{http://dx.doi.org/10.1006/jfan.1999.3456}{\textit{J.~Funct. Anal.}} \textbf{167} (1999), 243--344.

\end{thebibliography}
\end{document}
